\documentclass[12pt,a4paper]{article}
\usepackage[margin=1in]{geometry}
\usepackage{amsmath}
\usepackage{booktabs}
\usepackage{setspace}
\setstretch{1.18}

\title{WHA Core-1 (Value-Momentum Daily)\\Strategy Overview}
\author{}
\date{February 18, 2026}

\begin{document}
\maketitle

\section{Core Idea}
WHA Core-1 is a daily cross-sectional stock scoring strategy with a value-dominant, momentum-assisted structure.

\[
\mathrm{Score}_{i,t}=0.9\times\mathrm{ValueScore}_{i,t}+0.1\times\mathrm{MomentumScore}_{i,t}
\]

where:
\begin{itemize}
  \item \texttt{ValueScore}: valuation attractiveness (primary driver)
  \item \texttt{MomentumScore}: medium-term trend strength (secondary driver)
\end{itemize}

\section{Factor Construction}
\subsection*{1) Value block}
Three valuation metrics are used:
\begin{itemize}
  \item \texttt{PE} (Price-to-Earnings)
  \item \texttt{P/FCF} (Price-to-Free-Cash-Flow)
  \item \texttt{EV/EBITDA} (Enterprise Value / EBITDA)
\end{itemize}

To align direction (higher = cheaper), multiples are transformed into yield-style variables:
\[
\mathrm{earnings\_yield}=\frac{1}{\mathrm{PE}},\quad
\mathrm{fcf\_yield}=\frac{1}{\mathrm{P/FCF}},\quad
\mathrm{ev\_ebitda\_yield}=\frac{1}{\mathrm{EV/EBITDA}}
\]

\subsection*{2) Momentum block}
\begin{itemize}
  \item Lookback: 126 trading days (about 6 months)
  \item Skip window: latest 21 trading days (about 1 month) to reduce short-term reversal noise
\end{itemize}

\subsection*{3) Post-combination risk controls}
\begin{itemize}
  \item Industry neutralization (reduce sector-level systematic drift)
  \item Size neutralization (reduce large-cap/small-cap style bias)
  \item Beta neutralization (reduce broad market direction exposure)
  \item Quantile winsorization (limit outlier dominance)
\end{itemize}

\section{Single-Factor Backtest Evidence (Core Pair)}
Backtest window and segmentation:
\begin{itemize}
  \item Time range: 2010-01-01 to 2026-01-28
  \item Segmentation: 2-year slices, 9 slices total
\end{itemize}

\subsection*{A. Stage 1 (v2.1)}
\begin{center}
\begin{tabular}{lccc}
\toprule
Factor & IC Mean & IC Std & Positive IC Ratio \\
\midrule
Value\_v2 & 0.047520 & 0.015569 & 88.89\% \\
Momentum\_v2 & 0.012868 & 0.022771 & 66.67\% \\
\bottomrule
\end{tabular}
\end{center}

\subsection*{B. Stage 2 Strict (core pair)}
\begin{center}
\begin{tabular}{lccc}
\toprule
Factor & IC Mean & IC Std & Positive IC Ratio \\
\midrule
Value\_v2 & 0.053457 & 0.021938 & 100.00\% \\
Momentum\_v2 & 0.014055 & 0.026392 & 75.00\% \\
\bottomrule
\end{tabular}
\end{center}

Conclusion: Value is the stronger/stabler component; Momentum adds complementary information, leading to the final 0.9/0.1 linear mix.

\section{Three-Layer Combo Validation}
Locked combo: \texttt{linear}, \texttt{value=0.90}, \texttt{momentum=0.10}.

\subsection*{Layer 1: Segmented validation}
\begin{itemize}
  \item Range: 2010-01-01 to 2026-01-28
  \item Method: 2-year segmentation
  \item Outcome: linear combo outperformed non-linear candidates (gated / two-stage)
\end{itemize}

\subsection*{Layer 2: Fixed train/test (run\_with\_config)}
\begin{itemize}
  \item Train: 2010-01-04 to 2017-12-31
  \item Test: 2018-01-01 to 2026-01-28
  \item Train IC (overall): \texttt{0.080637}
  \item Test IC (overall): \texttt{0.053038}
\end{itemize}

\subsection*{Layer 3: Walk-forward validation}
Settings: \texttt{train=3y, test=1y, REBALANCE\_MODE=None}.
\begin{itemize}
  \item Rolling windows: 2013--2025 (n=13)
  \item Base range: start-year=2010, end-year=2025
  \item test\_ic: mean=\texttt{0.057578}, std=\texttt{0.033470}, pos\_ratio=\texttt{1.0000}
  \item test\_ic\_overall: mean=\texttt{0.050814}, std=\texttt{0.032703}, pos\_ratio=\texttt{1.0000}
\end{itemize}

\section{Final Stress Test (Post-WF)}
Stress profile:
\begin{itemize}
  \item \texttt{COST\_MULTIPLIER=1.5}
  \item \texttt{MIN\_MARKET\_CAP=2e9}
  \item \texttt{MIN\_DOLLAR\_VOLUME=5e6}
\end{itemize}

Results (n=13):
\begin{itemize}
  \item test\_ic: mean=\texttt{0.053310}, std=\texttt{0.032486}, pos\_ratio=\texttt{1.0000}
  \item test\_ic\_overall: mean=\texttt{0.046618}, std=\texttt{0.032058}, pos\_ratio=\texttt{1.0000}
\end{itemize}

Conclusion: the signal remained positive under stricter cost and liquidity constraints.

\section{Execution Convention and Practical Boundaries}
\subsection*{1) Data and universe}
\begin{itemize}
  \item Price: adjusted daily close (adjClose)
  \item Valuation: FMP ratio data transformed into yield-style features
  \item Base investability filters: minimum market cap, minimum dollar volume, minimum price
\end{itemize}

\subsection*{2) Trading convention}
\begin{itemize}
  \item Signal frequency: daily
  \item Execution delay: signal at T is used for T+1 (execution\_delay=1)
  \item Default portfolio orientation: long-only ranking
\end{itemize}

\subsection*{3) Known limitations}
\begin{itemize}
  \item Metrics here are research/backtest diagnostics, not full real-money PnL statements
  \item Live deployment should monitor rolling 20/40/60-day excess stability
  \item Capacity and impact analysis is required before scaling capital
\end{itemize}

\section{Current Completion Status (Disclosure-ready)}
\begin{itemize}
  \item Completed: single-factor validation, three-layer combo validation, and stress testing (cost multiplier + stricter liquidity constraints)
  \item Partially completed: cost stress outputs are available, but a full capacity-tier and impact-cost quant table is not yet standardized
  \item In progress: standardized risk-scenario casebook and live monitoring dashboard (20/40/60-day excess and turnover)
\end{itemize}

\section{One-line Summary}
WHA Core-1 is a value-first, momentum-assisted daily ranking strategy with positive evidence across single-factor tests, three-layer combo validation, and final stress testing; it is suitable as a daily signal engine with controlled live-scale rollout.

\end{document}
